\documentclass[11pt]{article}
\usepackage{amsmath,amsthm,amssymb}
\usepackage[margin=1in]{geometry}
\usepackage{booktabs}
\usepackage{array}
\usepackage{stmaryrd}
\usepackage[colorlinks=true,linkcolor=blue,citecolor=blue,urlcolor=blue]{hyperref}

\newtheorem{theorem}{Theorem}
\newtheorem{proposition}[theorem]{Proposition}
\newtheorem{corollary}[theorem]{Corollary}
\newtheorem{lemma}[theorem]{Lemma}
\theoremstyle{definition}
\newtheorem{definition}[theorem]{Definition}
\newtheorem{remark}[theorem]{Remark}
\newtheorem{example}[theorem]{Example}
\newtheorem*{program}{The nerve-level program}

\newcommand{\Poly}{\mathsf{Poly}}
\newcommand{\Cat}{\mathsf{Cat}}
\newcommand{\Set}{\mathsf{Set}}
\newcommand{\Ab}{\mathsf{Ab}}
\newcommand{\Comon}{\mathsf{Comon}}
\newcommand{\Ob}{\mathrm{Ob}}
\newcommand{\Hom}{\mathrm{Hom}}
\newcommand{\Aut}{\mathrm{Aut}}
\newcommand{\tr}{\mathrm{tr}}
\newcommand{\N}{\mathbb{N}}
\newcommand{\Z}{\mathbb{Z}}
\newcommand{\F}{\mathbb{F}}
\newcommand{\id}{\mathrm{id}}
\newcommand{\dom}{\operatorname{dom}}
\newcommand{\cod}{\operatorname{cod}}
\newcommand{\sS}{\mathcal{S}}
\newcommand{\cC}{\mathcal{C}}
\newcommand{\cD}{\mathcal{D}}
\newcommand{\lto}{\!\lhd\!}
\newcommand{\Nerve}{\mathrm{N}}

\title{Composition is a nerve-level-two datum\\
\large Grading compositional invariants of directed containers by nerve level,
not cohomological degree --- a program, with proved instances}
\author{MacBeth\thanks{Kodamai. This is a positioning note for the Kodamai
AI-Mathematician programme; the proofs it organizes are established in the cited companion notes
and proof files, which live together with this note and the verification scripts at
\texttt{github.com/scotmacbeth/work-in-progress} (content commit recorded on this page). It
\emph{supersedes and corrects} the earlier note \emph{``Composition is a degree-two datum''}
(content commit \texttt{ae9d67a}), whose headline grading was off by one; see the Introduction.
Prepared as a framing appendix for the Strathclyde mini-course \emph{Tangent Categories and
Containers} (13--16 October 2026).}\\
\small Kodamai}
\date{8 October 2026\\[2pt]
\small content commit \texttt{242c7ea}}

\begin{document}
\maketitle

\begin{abstract}
A directed container is a small category, and a question about such a category can be a question
about its underlying shape, about a resource transported across it, or about how its morphisms
compose. We ask whether a numerical grade on an invariant records how much of composition the
invariant can see, and we argue that the right grade is the \emph{nerve level} the invariant reads,
not --- as an earlier note of ours claimed --- its cohomological degree. The two indices differ by
one: a datum on composable pairs lives at nerve level two, but its cohomology is first read by a
degree-\emph{one} cocycle condition. Out-degree homogeneity reads only the reflexive graph (nerve
level one) and is provably blind to composition. The comultiplication $\delta$ is a nerve-level-two
datum, and the invariants built from it --- the left-translation trivialization of the tangent
bundle, detecting groupoids and, one notch finer, thin groupoids --- separate the group $\Z/2$ from
the idempotent monoid $\{1,e\}$ at cohomological degree \emph{one}, exactly the separation the floor
cannot make. The Zappa--Sz\'ep weld $[\omega]\in H^2$ is a nerve-level-\emph{three} datum: it reads
associativity, a condition on composable triples. We prove that the truncation tower collapses at
level two (a category is its $2$-truncated nerve), so above the floor there is no third
\emph{detection} rung; the genuine level-two/level-three distinction is the condition-level fact that
building a category from a composition is an associativity problem. Finally we show the degree-one
inventory class $[\theta_R]$ and the degree-two weld class $[\omega]$ are two independent anchors in
the composition-sensitive band, differing in coefficients and condition-level, not in whether they
see composition. We assemble these as a program --- nerve level grades compositional data --- and
flag, as an open question with a degree-one counterexample, whether cohomological \emph{degree}
stratifies \emph{detection} above the floor.
\end{abstract}

\section{Introduction}\label{sec:intro}

\paragraph{A correction, stated first.}
An earlier note of ours carried the headline \emph{``Composition is a degree-two datum,''} and with
it the claim that \emph{no invariant below cohomological degree two can detect composition}. The
claim is false, and this note replaces it. The note's own degree-\emph{one} invariant already
separates the group $\Z/2$ from the idempotent monoid $\{1,e\}$ --- a pair that differs only in how
its one nontrivial arrow composes with itself. The error was a conflation of two distinct simplicial
indices attached to a cohomology class, and the fix is to grade by the right one. The underlying
theorems are unaffected; only the grading, and the slogan that named it, were wrong. We state the
corrected grading, prove the regrade that forces it, and rebuild the ladder honestly.

\paragraph{One functor, three questions.}
A directed container is, by the theorem of Ahman and Uustalu \cite{AhmanUustalu2016}, exactly a
small category: the comonoids for the composition product $\lhd$ on polynomial functors are small
categories, with objects the shapes and, for each object $a$, the direction set $E_a$ as the set of
morphisms out of $a$. A single such category carries structure at several levels at once. There is
its \emph{underlying functor} $p=\sum_a y^{E_a}$, which remembers only the object-indexed family of
direction sets. There is \emph{data transported across it} --- a presheaf of resources, say stock
levels at the nodes of a supply chain. And there is its \emph{composition} --- the comultiplication
$\delta\colon p\to p\lto p$, which records how two morphisms compose, and, for a composite built as a
Zappa--Sz\'ep product, the weld that glues two factors together. Questions can interrogate any of
these levels, and a recurring experience of this programme is that an invariant designed for one is
\emph{structurally unable} to answer another. The derivative is blind to composition; the
consistency of a transported resource is independent of whether the composite factors as claimed.
These are facts about where the relevant datum lives.

\paragraph{The two indices, and the regrade.}
Make that experience precise by grading an invariant by the \emph{simplicial level of the nerve} the
datum it reads occupies. The nerve $\Nerve(\cC)$ of a category has $\Nerve_0$ the objects,
$\Nerve_1$ the arrows, and $\Nerve_k$ the composable $k$-tuples; composition is the inner face
$d_1\colon\Nerve_2\to\Nerve_1$, so the comultiplication $\delta$ is a datum on $\Nerve_2$, the
composable pairs. This is the ``two'' that the old slogan heard. But the \emph{cohomological degree}
of a class is a different index, and it is lower by one:

\begin{quote}
\emph{An $n$-cochain is a function on $\Nerve_n$, but the cocycle condition $d^n c=0$ that defines its
class is a system of equations indexed by $\Nerve_{n+1}$. The binding simplicial level of an
$H^n$-class is thus $\Nerve_{n+1}$: cohomological degree $=$ (nerve level of the cocycle condition)
$-\,1$.}
\end{quote}

Composable pairs ($\Nerve_2$) are therefore read already by a degree-\emph{one} cocycle condition.
The old ``degree two'' took the nerve level of $\delta$ and mislabelled it as a cohomological degree.
We correct this throughout (\S\ref{sec:prelim}, Lemma~\ref{lem:regrade}); the retired slogan becomes:

\begin{quote}
\emph{Composition --- the datum of composable pairs --- is a nerve-level-two datum, first detectable
at cohomological degree one. Associativity --- the datum of composable triples --- is the
nerve-level-three datum, cohomological degree two, and the Zappa--Sz\'ep weld
$[\omega]\in H^2$ is its obstruction.}
\end{quote}

\paragraph{The honest ladder.}
The detection ladder has, above the floor, \emph{two} rungs for the price the old note charged for
three, and a genuine third rung that is not a detection rung at all.
\begin{itemize}
\item \textbf{Nerve level $1$ (the blind floor).} Invariants that factor through the reflexive graph
$\tr_1\Nerve(\cC)$ --- equivalently through the forgetful $U\colon\Cat\to\Poly$ --- read out-degrees
and nothing about composition. Tangent-regularity lives here and is provably blind: $\Z/2$ and
$\{1,e\}$ have the same underlying functor $y^2$ and are indistinguishable to it.
\item \textbf{Nerve level $2$ (composition).} Invariants that read the composition map
$d_1\colon\Nerve_2\to\Nerve_1$ can see composition. Two live here. The $\delta$-translation
trivialization detects groupoids, and one notch finer thin groupoids, and it \emph{separates $\Z/2$
from $\{1,e\}$ --- at cohomological degree one.} The inventory class $[\theta_R]\in H^1(\sS;A)$ of a
transported resource is the other.
\item \textbf{Nerve level $3$ (associativity).} The weld $[\omega]\in H^2$ reads a condition on
composable triples: whether a composition assembles into an associative category. For a \emph{fixed}
category this rung is empty of new detection power --- a category is reconstructible from its
$2$-truncated nerve (Theorem~\ref{thm:coskel}), so the truncation tower collapses at level two.
The genuine level-two/level-three distinction is in the \emph{assembly} problem, where the object
varies: composition exists ($\Nerve_2$) versus composition is associative ($\Nerve_3$).
\end{itemize}

\paragraph{This is a program, not a theorem.}
What is proved is a regrade lemma, a strict floor separation, a collapse theorem, an assembly
separation, three invariant instances, and one independence result --- each in a companion note or
proof file cited below and not re-derived here. What is \emph{not} proved, and what we are careful to
state as an open question (\S\ref{sec:program}), is that cohomological \emph{degree} stratifies
\emph{detection} power above the floor. As stated naively (``higher degree detects strictly more'')
this is \textbf{false}: the degree-one $\delta$-invariant already detects composition, and the
degree-one and degree-two cohomological anchors are \emph{independent}, not ordered. The grade that
does the stratifying is \emph{nerve level}; cohomological degree is a coefficient-dependent shadow of
it (Remark~\ref{rem:coefftest}). This is a positioning note: its job is to make the pattern visible
and correctly labelled, not to establish new mathematics.

\paragraph{Outline.} Section~\ref{sec:prelim} fixes containers as categories, the nerve and its two
truncations, the regrade lemma, the two structure maps $\partial$ and $\delta$, the resource torsor
and its $H^1$, the weld and its $H^2$, and the nerve-level grading.
Section~\ref{sec:lvl1} is the blind floor; Section~\ref{sec:lvl2} is composition, where the
$\delta$-translations separate what the floor cannot and the monodromy tower stratifies them;
Section~\ref{sec:lvl3} is associativity, the weld, and the collapse/assembly facts;
Section~\ref{sec:indep} shows the degree-one inventory class and the degree-two weld class are
independent. Section~\ref{sec:program} states the program and the open question.

\section{Preliminaries}\label{sec:prelim}

We work with polynomial functors in one variable and finite support; everything is first order.

\paragraph{Containers as categories.}
A \emph{container} in one variable is a polynomial functor $p=\sum_{a\in\Ob}y^{E_a}$ with a set
$\Ob$ of \emph{shapes} and, for each shape $a$, a set $E_a$ of \emph{directions}
\cite{AAGM2003,GambinoKock2013}. The composition product $\lhd$ has unit $y$ and realizes
substitution. By Ahman and Uustalu \cite{AhmanUustalu2016} the comonoids in $(\Poly,\lhd,y)$ are
exactly small categories: a comonoid structure $\delta\colon p\to p\lto p$ chooses identities (the
counit) and a composition (the comultiplication), $\Ob$ is the object set, and
$E_a=\coprod_c\Hom(a,c)$ is the set of morphisms out of $a$, so $|E_a|$ is the \emph{out-degree} of
$a$. Write $\Comon(\Poly,\lhd)\simeq\Cat$ for this equivalence and $U\colon\Cat\to\Poly$ for the
forgetful functor that discards $\delta$ and keeps only $\sum_a y^{E_a}$. We compose
\emph{diagrammatically}: for $g\colon a\to b$ and $h\colon b\to c$, the composite $g\cdot h\colon
a\to c$ means ``$g$ then $h$''.

\paragraph{The nerve and its two truncations.}
The \emph{nerve} $\Nerve(\cC)$ is the simplicial set with $\Nerve_0=\Ob$, $\Nerve_1$ the arrows, and
$\Nerve_k$ the composable $k$-tuples $a_0\xrightarrow{f_1}\cdots\xrightarrow{f_k}a_k$. The inner face
$d_1\colon\Nerve_2\to\Nerve_1$ composes the pair, $(f,g)\mapsto f\cdot g$; the degeneracies insert
identities. The \emph{$k$-truncation} $\tr_k\Nerve(\cC)$ keeps $\Nerve_0,\dots,\Nerve_k$ with all
faces and degeneracies among them. Two truncations matter here:
\[
   \tr_1\Nerve(\cC)=\text{the reflexive graph of }\cC\ (\text{objects, arrows, source, target,
   identities})=U(p);
\]
$\tr_2\Nerve(\cC)$ adds the composition map $d_1$. For an abelian coefficient system $A$ (a functor
$\cC\to\Ab$; constant unless noted), $H^n(\cC;A)$ is the cohomology of the classifying space $B\cC$,
computed by the normalized cochain complex with $C^n=\{$functions on nondegenerate $n$-simplices$\}$
and the standard bar differential; for a one-object category it is monoid/group cohomology.

\paragraph{The regrade lemma.}
The grade we use is the nerve level the \emph{binding} datum of an invariant occupies. Cohomological
degree is a strictly lower index, and the relation is exact.

\begin{lemma}[regrade]\label{lem:regrade}
For every $n\ge 0$, the cocycle condition $d^n c=0$ defining an $H^n$-class is a system of equations
indexed by $\Nerve_{n+1}$. Consequently
\[
   \text{cohomological degree }n \;=\; (\text{nerve level of the cocycle condition})\;-\;1.
\]
In particular a datum on composable pairs ($\Nerve_2$) is read by a degree-\emph{one} cocycle
condition, and one on composable triples ($\Nerve_3$) by a degree-two cocycle condition.
\end{lemma}

\begin{proof}[Reason]
$d^n\colon C^n\to C^{n+1}$ sends a function on $\Nerve_n$ to a function on $\Nerve_{n+1}$: by the bar
formula $(d^n c)(f_1,\dots,f_{n+1})$ is an alternating sum of $c$ on the $n$-faces of the
$(n+1)$-simplex $(f_1,\dots,f_{n+1})$, together with the coefficient action along $f_1$. Hence
$d^n c=0$ is the family $\{(d^n c)(\sigma)=0:\sigma\in\Nerve_{n+1}\}$, indexed by $\Nerve_{n+1}$; the
coboundary relation reads only $\Nerve_n$, so the highest level entering the definition of the class
is $\Nerve_{n+1}$. Proved in full in \cite{MacBethNerve}.
\end{proof}

\begin{remark}[the ``two'' of $p\lto p$]\label{rem:simpvsdeg}
The comultiplication $\delta\colon p\to p\lto p$ lands in the two-fold composite $p\lto p$: it is data
on composable \emph{pairs}, i.e.\ on the $2$-simplices of the nerve, and in that \emph{simplicial}
sense composition is degree-two data. Cohomological degree is the other index: by
Lemma~\ref{lem:regrade} an $H^1$ cocycle condition already reads $2$-simplices. The earlier note
named the simplicial level ``degree two'' and thereby read $\delta$ one cohomological degree too high;
every ``degree two'' of that note that referred to composition itself should read ``nerve level two,
cohomological degree one''.
\end{remark}

\paragraph{Two structure maps: $\partial$ reads the graph, $\delta$ reads composition.}
The one-hole derivative \cite{AAGM2003} is
\[
   \partial p \;=\; \sum_a |E_a|\cdot y^{\,|E_a|-1},
\]
depending only on the direction sets, i.e.\ only on $U(p)=\tr_1\Nerve(\cC)$: a construction on the
reflexive graph, nerve level one. With $W=\Z[\varepsilon]/(\varepsilon^2)$ the dual-number Weil
algebra, the attached tangent functor is the Weil base change $T=(-)\otimes W$; the one-hole
derivative appears as the $\varepsilon$-coefficient of the substitution $p(y+\varepsilon v)$, and all
we use downstream is that the vertical fibre over a shape $a$ is free of rank $|E_a|$
\cite{MacBethCartan,MacBethParallel}. By contrast $\delta\colon p\to p\lto p$ is \emph{not} a
construction on $U(p)$: it is the composition map $d_1$, a datum on $\Nerve_2$. This is the precise
sense in which composition is a nerve-level-two datum, and the datum $\partial$ cannot see.

\paragraph{The resource torsor and its degree-one class.}
Model a supply chain as a small category $\sS$: objects are nodes, morphisms out of $a$ are
admissible operations, composition is sequencing. \emph{Inventory} is data transported across $\sS$:
a local system in which each operation $f$ acts on stock by a translation $R(f)(x)=x+\tau_f$ of an
abelian stock-group $A$. (We write $\tau_f$ for the translation $1$-cochain to avoid any clash with
the comultiplication $\delta$.) A \emph{global section} is a node-potential $\varphi$ with
\[
   \tau_f \;=\; \varphi(\dom f)-\varphi(\cod f)\qquad\text{for every }f,
\]
i.e.\ the $1$-cochain $\tau$ is a coboundary. Because the translations are not group homomorphisms,
the fibrewise data form a \emph{torsor} under the linear local system $A$, not a presheaf of groups;
the obstruction is accordingly the class
\[
   [\theta_R]\;:=\;[\tau]\in H^1(\sS;A).
\]
This is a degree-one class, coefficient system the stock-group $A$, measuring route/reconvergence
discrepancy over the nerve of $\sS$ --- the classical mechanism of Kirchhoff's voltage law
(\S\ref{sec:indep}).

\paragraph{The weld and its degree-two class.}
A composite chain is a \emph{Zappa--Sz\'ep product} $\sS=\cC\bowtie\cD$ (procurement $\bowtie$
logistics): $\cC,\cD$ are wide subcategories meeting in the discrete category on the shared objects,
and every morphism factors uniquely as $f=c\cdot d$ \cite{AhmanUustalu2016,MacBethZS}. Whether a
candidate matched pair of actions closes into such a factorization is governed by the \emph{weld
obstruction}
\[
   [\omega]\in H^2(\sS;\cD),
\]
a Baues--Wirsching $2$-cocycle on the orbit category with coefficients the vertex-group system $\cD$;
the chain decomposes as claimed exactly when $[\omega]=0$ \cite{MacBethZS}. By
Lemma~\ref{lem:regrade} this degree-two class reads $\Nerve_3$: its cocycle condition \emph{is} the
associativity equation of the assembled composition (\S\ref{sec:lvl3}). It is the home of the
blockchain re-entrancy obstruction $[\omega]=\varepsilon\in H^2\cong\Z/2$ \cite{MacBethReentrancy}.

\paragraph{The nerve-level grading.}
We rank a compositional invariant by the nerve level of the binding datum it reads: level $1$ for the
reflexive graph $U(p)$ (out-degree family); level $2$ for the composition map $d_1$ --- read by the
comultiplication $\delta$ and by any $H^1(\sS;A)$ of a transported local system; level $3$ for the
associativity of the assembled composition, whose obstruction is the weld $[\omega]\in H^2$. The grade
is a grade of \emph{data}; the different cohomological degrees carry different coefficient systems, and
as Remark~\ref{rem:nocollapse} records, the same cyclic group can surface at different degrees under a
change of coefficients, so one must not read the grading as a single filtration.

\section{Nerve level one: the invariant blind to composition}\label{sec:lvl1}

The floor reads only the reflexive graph and is for that reason unable to see composition at all. The
result is proved in \cite{MacBethParallel}; we state it and exhibit the discriminating examples.

The tangent bundle $T p=p\otimes W$ has, over a shape $a$, a fibre free of rank $|E_a|$. Call $p$
\emph{tangent-regular} if these fibres are all isomorphic --- the honest residue, in the
representable/monoidal setting, of parallelizability in the sense of Cruttwell, Ikonicoff, Lemay and
Van Der Linden \cite[Def.~4.18, attributed there]{Lanfranchi2026}.

\begin{theorem}[\cite{MacBethParallel}]\label{thm:lvl1}
A container $p=\sum_a y^{E_a}$ is tangent-regular iff the out-degree $a\mapsto|E_a|$ is constant on
$\Ob$. This condition depends only on $U(p)=\tr_1\Nerve(\cC)$ --- it factors through
$U\colon\Cat\to\Poly$ --- and is therefore blind to composition: it is invariant under changing
$\delta$ on a fixed underlying functor $p$, and so cannot be equivalent to any property that
distinguishes two comonoid structures on the same $p$, such as being a groupoid
(Example~\ref{ex:lvl1}).
\end{theorem}

\begin{example}[the blindness, made concrete]\label{ex:lvl1}
The idempotent monoid $\{1,e\}$ with $e^2=e$ and the group $\Z/2$ both have underlying functor $y^2$.
They have the \emph{same} tangent bundle and are both tangent-regular (one object, so out-degree is
trivially constant), yet $\Z/2$ is a groupoid and $\{1,e\}$ is not. The floor cannot tell them apart.
Conversely the groupoid $\Z/2\sqcup\Z/3$ has out-degrees $2$ and $3$, so $U=y^2+y^3$ is \emph{not}
tangent-regular though it is a groupoid: tangent-regularity neither implies nor is implied by being a
groupoid.
\end{example}

The cautionary form of this rung is easy to fall for. A genuine theorem of Lanfranchi
\cite[Thm.~4.20]{Lanfranchi2026} --- \emph{for a group object $G$ in a Cartesian tangent category, $G$
is a Lie group object iff $G$ is parallelizable} --- invites the port ``a container is parallelizable
iff it is a groupoid.'' The hypothesis ``given a group object'' is essential: without it the statement
is false ($S^7$ is parallelizable with no group structure), and the port drops it. The port is in any
case ill-posed, because the container derivative's tangent structure is the monoidal base change
$(-)\otimes W$ and not a Cartesian one \cite{MacBethParallel}; what survives is tangent-regularity,
which lives at nerve level one and cannot see the groupoid property. This is no defect in
\cite{Lanfranchi2026} --- that theorem is about group objects in Cartesian tangent categories --- but
the symptom of asking a graph-level invariant a composition-level question. Connected groupoids
\emph{are} tangent-regular; but strictly more categories are, and the gap is the composition content
that only nerve level two recovers (\S\ref{sec:lvl2}).

\section{Nerve level two: composition, and the invariant that sees it --- at degree one}
\label{sec:lvl2}

The second rung reads the composition map $d_1$. Here composition becomes visible, and --- this is the
crux of the correction --- it becomes visible already at cohomological degree \emph{one}. We first
exhibit the $\delta$-translation invariant that separates the pair the floor could not, then verify by
direct computation that the separation is a degree-one cohomological fact, and finally refine the
invariant into a monodromy tower.

\paragraph{The $\delta$-translation invariant.}
Fix a morphism $g\colon a\to b$ and pin the first slot of $\delta$ to it: the
\emph{$\delta$-translation}
\[
   L_g\colon E_b\to E_a,\qquad L_g(h)=g\cdot h,
\]
precomposition by $g$. It is contravariantly functorial, assembling $\delta$ into the
\emph{arrows-out presheaf} $E_{(-)}\colon\cC^{\mathrm{op}}\to\Set$, $a\mapsto E_a$, $g\mapsto L_g$ --- a
property-free restatement of the composition map $d_1$. (The statement below concerns this
vertical/linear layer of $T p=p\otimes W$, not the full bundle \cite{MacBethDelta}.) Say $Tp$ is
\emph{$\delta$-trivializable} over a connected component $[a]$ if $E_{(-)}$ restricted there is a
\emph{local system} --- every $L_g$ a bijection --- the honest discrete analogue of trivializing a
tangent bundle by translations, $TG\cong G\times\mathfrak g$.

\begin{theorem}[\cite{MacBethDelta}]\label{thm:lvl2}
$Tp$ is $\delta$-trivializable over $[a]$ iff the full subcategory on $[a]$ is a \emph{groupoid};
globally, iff $\cC$ is a groupoid. Surjectivity of the translations already forces full invertibility:
a surjective $L_g$ hits $\id_a\in E_a$, so some $h$ has $g\cdot h=\id_a$; surjectivity of $L_h$ then
gives $k$ with $h\cdot k=\id_b$, whence $g=g\cdot(h\cdot k)=(g\cdot h)\cdot k=k$ and so $h\cdot g=\id_b$
--- $g$ is invertible. The dividing line is exactly ``groupoid'', with no weaker cancellation
condition between.
\end{theorem}

\begin{example}[seeing what the floor could not]\label{ex:lvl2}
Return to the pair of Example~\ref{ex:lvl1}. On $\Z/2$ every $L_g$ is a bijection, so $Tp$ is
$\delta$-trivializable. On $\{1,e\}$ the translation $L_e$ sends $1\mapsto e$ and $e\mapsto e$, missing
$\id=1$; it is not even injective, and $\{1,e\}$ is not $\delta$-trivializable. The invariant
\emph{separates $\Z/2$ from $\{1,e\}$} --- the exact pair the floor, reading only the common underlying
functor $y^2$, was structurally unable to tell apart (Theorem~\ref{thm:lvl1}). This is the content of
the slogan we can prove: there is a genuine invariant, built from $\delta$ (the composition map $d_1$),
that detects a composition property no invariant factoring through $U$ can.
\end{example}

\paragraph{The separation is a degree-one fact.}
It remains to pin the cohomological degree at which this separation occurs, and to show it is
\emph{one}, not two. This is Theorem~2 of \cite{MacBethNerve}; we record the computation because it is
where the old slogan broke.

\begin{theorem}[\cite{MacBethNerve}]\label{thm:degreeone}
The categories $\Z/2$ and $\{1,e\}$ have isomorphic $1$-truncated nerves (hence agree on every
invariant factoring through $U$, in particular on out-degree), but
$\dim_{\F_2}H^1(-;\F_2)$ distinguishes them: $H^1(\Z/2;\F_2)=\F_2$ while $H^1(\{1,e\};\F_2)=0$. The
distinction is a function of $\tr_2\Nerve$ and not of $\tr_1\Nerve$. Hence detection of composition
begins at nerve level two, which by Lemma~\ref{lem:regrade} is cohomological degree \emph{one}.
\end{theorem}

\begin{proof}[Computation]
The unique nondegenerate composable pair is $(a,a)$ with $d_1(a,a)=a\cdot a$, which is $1$ in $\Z/2$
and $a$ in $\{1,e\}$; so no isomorphism of $\tr_2$ nerves exists, while the relabelling of
$\{1,a\}$ is an isomorphism of $\tr_1$ nerves. Over $\F_2$ a normalized $1$-cochain is a value $c(a)$.
For $\{1,e\}$: $(d^1 c)(e,e)=c(e)-c(e\cdot e)+c(e)=c(e)$, forcing $c(e)=0$, so $H^1=0$. For $\Z/2$:
$(d^1 c)(a,a)=2c(a)-c(1)=0$ over $\F_2$, so every value is a cocycle and $H^1=\F_2$. Full details and
a verifying script in \cite{MacBethNerve}.
\end{proof}

\begin{remark}[why \emph{nerve level}, not degree, is the grade --- the coefficient test]
\label{rem:coefftest}
The degree at which this single $\Nerve_2$ distinction surfaces \emph{depends on the coefficients}.
The feature $a\cdot a=1$ creates a $1$-cycle, $H_1(\Z/2)=\Z/2$, living at nerve level $\Nerve_2$. Over
$\F_2$ it is detected by $H^1$ (degree one); over $\Z$ the same torsion is detected by $H^2$ (degree
two, $H^2(\Z/2;\Z)=\Z/2$), the Bockstein having moved it up a degree. So ``which degree detects
composition here'' has no coefficient-free answer --- but ``which \emph{nerve level}'' always does:
$\Nerve_2$. This is the thesis in one example, and it is also why the old ``degree two'' was not even
stably wrong: over $\F_2$ the detection is degree one.
\end{remark}

\paragraph{A monodromy tower on the $\delta$-translations.}
The $\delta$-translations can be asked to do more than be invertible: to trivialize to a
\emph{constant} identification, with trivial monodromy. Call $E_{(-)}$
\emph{constant-trivializable} over $[a]$ if it is naturally isomorphic to a constant presheaf
$\Delta_V$: bijections $\varphi_b\colon E_b\to V$ with $\varphi_x\circ L_g=\varphi_y$ for \emph{every}
$g\colon x\to y$, loops included. Read against a loop this asks the translation cocycle $\{L_g\}$ to
be a coboundary.

\begin{theorem}[\cite{MacBethThin}]\label{thm:thin}
$E_{(-)}$ is constant-trivializable over $[a]$ iff $[a]$ is a \emph{thin} (indiscrete) groupoid ---
trivial vertex groups, i.e.\ an equivalence-relation component; globally, iff $\cC$ is a thin
groupoid. Constant-trivializability implies $\delta$-trivializability, so the groupoid hypothesis of
Theorem~\ref{thm:lvl2} comes for free; applying naturality to a loop $g\in\Aut(a)$ forces $L_g=\id$,
and since the regular action fixes the basepoint $\id_a\in E_a$, forces $g=\id_a$ --- the vertex group
vanishes.
\end{theorem}

\begin{corollary}[the monodromy tower]\label{cor:tower}
Over a connected component, three successively finer conditions on the $\delta$-translation presheaf
$E_{(-)}$ stratify the composition datum:
\begin{center}
\begin{tabular}{@{}clll@{}}
\toprule
Rung & Condition on $E_{(-)}\big|_{[a]}$ & Equivalent to & Status \\
\midrule
$1$ & generic presheaf (each $L_g$ a function) & nothing & trivial \\
$2$ & local system (every $L_g$ a bijection) & $[a]$ a groupoid & \cite{MacBethDelta} \\
$3$ & $\cong$ constant presheaf & $[a]$ a thin groupoid & \cite{MacBethThin} \\
\bottomrule
\end{tabular}
\end{center}
Rungs $2$ and $3$ are distinct, separated by a nontrivial vertex group: $\Z/2$ is a groupoid but not
constant-trivializable (the regular action $\Z/2\curvearrowright\Z/2$ is free, nontrivial loop
monodromy), while the indiscrete two-object category $I_2$ is both. The obstruction to upgrading rung
$2$ to rung $3$ is the \emph{class of the regular representation} in the non-abelian
$H^1\!\big([a];\Sigma_{E_a}\big)$; this class is $\Hom(\Aut(a),\Sigma_{E_a})/\text{conj}$ evaluated at
the regular representation, which is faithful, so it vanishes iff $\Aut(a)=1$. Thus $\Aut(a)$ is the
\emph{holonomy group} of the tower, and it is a nontrivial \emph{vertex group} --- not
one-object-ness --- that marks the step.
\end{corollary}

So the composition datum carries an internal grading of its own: generic $\subsetneq$ groupoid
$\subsetneq$ thin groupoid, all three living at nerve level two as refinements of the
$\delta$-translation presheaf. Rung $3$ has no classical counterpart among positive-dimensional Lie
groups --- every Lie group already admits a global flat left-invariant framing with trivial holonomy
--- because in the discrete setting rung $3$ forces the group acting on the fibre to be trivial. We do
not overclaim: rungs $2$ and $3$ are composition-sensitive invariants of $\delta$, and the literal
cohomological anchor of rung $3$ is the non-abelian degree-one class above, not a new $H^2$ class. The
literal abelian anchors of the ladder remain $[\theta_R]\in H^1$ and $[\omega]\in H^2$; the tower
refines the \emph{datum} $\delta$ at nerve level two.

\section{Nerve level three: associativity is the weld}\label{sec:lvl3}

The weld $[\omega]\in H^2$ reads one level higher. By Lemma~\ref{lem:regrade} its degree-two cocycle
condition is a system of equations on composable triples $\Nerve_3$, and that system is exactly
associativity. We make the level-two/level-three distinction precise, and show why it is \emph{not} a
third detection rung for a fixed category.

\begin{theorem}[2-coskeletal collapse, \cite{MacBethNerve}]\label{thm:coskel}
A small category is determined up to isomorphism by its $2$-truncated nerve. Consequently, for
isomorphism-invariants of a fixed category, every invariant factors through $\tr_2\Nerve$: there is no
invariant that needs $\Nerve_3$, and the truncation tower stops at level two.
\end{theorem}

\begin{proof}[Reason]
From $\tr_2\Nerve(\cC)$ one recovers objects, arrows, source, target, identities, and composition
$d_1\colon\Nerve_2\to\Nerve_1$; these are all the data of a category, and an isomorphism of
$2$-truncated nerves is exactly an isomorphism of categories. Higher $\Nerve_k$ are forced by the
Segal condition. See \cite{MacBethNerve}.
\end{proof}

So the third rung is not a detection separation of \emph{fixed} categories. It is real, but it lives
in the \emph{assembly} problem, where the object varies: one is handed a composition on composable
pairs and asks whether it assembles into a category.

\begin{theorem}[associativity is an $\Nerve_3$-predicate, \cite{MacBethNerve}]\label{thm:assembly}
There is a unital magma whose full composition on composable pairs is defined and unital, yet which is
not associative --- e.g.\ on $\{1,x,y\}$ with unit $1$ and $x\cdot x=y$, $x\cdot y=x$, $y\cdot x=y$,
$y\cdot y=y$: the triple $(x,x,x)$ gives $(x\cdot x)\cdot x=y$ but $x\cdot(x\cdot x)=x$. Every pair-level
($\Nerve_2$) datum is present and unital; the obstruction lives purely at $\Nerve_3$. When the
composition arises as a Zappa--Sz\'ep weld, this obstruction is the class $[\omega]\in H^2$, vanishing
iff the weld is associative (a category). The separation of ``composition exists'' ($\Nerve_2$) from
``composition is associative'' ($\Nerve_3$) is therefore strict, witnessed by the re-entrancy weld
$[\omega]=\varepsilon\neq 0\in H^2\cong\Z/2$ \cite{MacBethReentrancy}.
\end{theorem}

This is the honest content of ``nerve level two $\subsetneq$ nerve level three'': not that a finer
invariant of a fixed category needs triples --- Theorem~\ref{thm:coskel} forbids that --- but that the
\emph{obstruction to building} a category from a composition is an associativity datum one level above
the composition itself. It is on this reading, and only this one, that the weld deserves to be called
the level-three (degree-two) datum of the programme.

\section{The two cohomological anchors are independent}\label{sec:indep}

The ladder has two literal abelian cohomological anchors: the inventory class
$[\theta_R]\in H^1(\sS;A)$ at nerve level two (a transported-resource datum) and the weld class
$[\omega]\in H^2(\sS;\cD)$ at nerve level three (an associativity datum). They are not ordered by
detection power; they are \emph{independent}. The results are proved in \cite{MacBethSupplyChain}.

\begin{theorem}[\cite{MacBethSupplyChain}]\label{thm:consistency}
Let $\sS=\cC\bowtie\cD$ and let $R$ be a resource local system with abelian stock-group $A$. Global
inventory is consistent --- $R$ admits a global section over the cover $\{\cC,\cD\}$ --- iff
$[\theta_R]=0\in H^1(\sS;A)$. The matched-pair law plays no role: a family compatible with the
$\cC$- and $\cD$-morphisms is a global section, because compatibility is closed under composition and
$\cC\cup\cD$ generates $\sS$.
\end{theorem}

\begin{theorem}[independence, \cite{MacBethSupplyChain}]\label{thm:sep}
$[\theta_R]\in H^1(\sS;A)$ and $[\omega]\in H^2(\sS;\cD)$ are logically independent: both
``$[\omega]=0$ with $[\theta_R]\neq 0$'' and ``$[\omega]\neq 0$ with $[\theta_R]=0$'' occur. Moreover
no family of homomorphisms carries either class to the other: by the two witnesses, either direction
would have to send a zero class to a nonzero class.
\end{theorem}

\begin{example}[the two witnesses]\label{ex:indep}
For the first direction, let $\sS$ be the free category on the diamond quiver
$a\rightrightarrows\{b,c\}\rightrightarrows d$ with distinct routes: directed-acyclic, so no
non-identity invertibles, the weld coefficient system is trivial and $[\omega]=0$ is forced; yet
$B\sS$ has first Betti number $1$, so a transport system with nonzero loop holonomy gives
$[\theta_R]\neq 0$. The second direction is \emph{generic}: for \emph{any} $\sS$ with $[\omega]\neq 0$,
the zero inventory system has $[\theta_R]=0$. A concrete such $\sS$ is the \emph{rigid-twist category}
--- one object with endomorphism monoid a copy of $\Z/2$ generated by a self-inverse twist --- whose
weld class generates $H^2(\sS;\Z/2)$. The substantive direction is thus the first; the second needs
only the freedom to choose the resource.
\end{example}

\begin{remark}[different degrees do not by themselves mean independence]\label{rem:nocollapse}
It would be a mistake to argue independence from the bare fact that the two classes sit in different
degrees. Standard degree-raising maps exist --- the transgression $d_2\colon H^1\to H^2$ of a
Lyndon--Hochschild--Serre spectral sequence, and Bockstein homomorphisms $H^1\to H^2$ --- so a priori
a class could be carried up a degree. Independence here is the \emph{content of the two witnesses}
(Theorem~\ref{thm:sep}), not a formal consequence of the degree gap. What the degree gap does explain
is why no \emph{isomorphism} could identify the two classes outright: for the free loop $\sS=B\N$ (one
loop; $B\sS\simeq S^1$) one has $H^1(B\N;A)=A$ while $H^2(B\N;A)=0$, and for a weld coefficient one
takes $\cD=\Z$ with $H^2(B\N;\Z)=0$ --- the degrees genuinely do not match. This is the same ``rhyme,
not identity'' caution that governs the two $H^2$-sites elsewhere in the programme \cite{MacBethZS},
here across one degree: the same cyclic group can appear on both rungs, but the site, the coefficients,
and the degree differ, and one must check all three before identifying two obstructions.
\end{remark}

The conjecture tested in \cite{MacBethSupplyChain} --- that inventory consistency \emph{equals} the
weld $[\omega]=0$ --- was a false generalization of blockchain re-entrancy, where the conserved
quantity \emph{is} the weld-state and no separate resource presheaf exists. In a supply chain the
stock lives on the nodes, transported by operations, and is a genuinely separate degree-one datum. The
modelling is sharpened, not debunked: re-entrancy is a nerve-level-three failure (handoff order, in
the weld), inventory inconsistency a nerve-level-two failure (route discrepancy, on the nodes) --- two
economic failure modes at two levels, each with a witness.

\section{The program, and what it would take to prove}\label{sec:program}

\paragraph{What is assembled here.}
Three invariant instances fall on a single nerve-level ladder: out-degree homogeneity at level one,
blind to composition; the $\delta$-translation invariant at level two, which detects groupoids and
thin groupoids and separates $\Z/2$ from $\{1,e\}$ --- at cohomological degree one; and the weld at
level three, the associativity obstruction. Two non-collapse facts make the ladder informative rather
than decorative: level-one blindness (Theorem~\ref{thm:lvl1}), and the independence of the two
cohomological anchors (Theorem~\ref{thm:sep}). A third, collapse, keeps us honest:
Theorem~\ref{thm:coskel} shows there is no third \emph{detection} rung for a fixed category.

\paragraph{The program, stated as a program.}
\begin{program}
For a directed container (small category) $\cC$ and a compositional invariant $I$, the \emph{nerve
level} of the binding datum $I$ reads records what $I$ can detect: an invariant factoring through the
$k$-truncation $\tr_k\Nerve(\cC)$ is blind to every composition phenomenon genuinely of nerve level
$>k$. In particular composition --- the comultiplication $\delta$, the composition map $d_1$ --- is a
nerve-level-two datum, first detectable at cohomological degree one; and associativity --- the weld
$[\omega]\in H^2$ --- is the nerve-level-three datum, cohomological degree two.
\end{program}
This is a \textbf{conjecture and research direction, not a theorem}. We have a regrade lemma, a strict
floor separation, a collapse theorem, an assembly separation, three instances and one independence
fact; we do not have a uniform ``blindness'' theorem that an invariant factoring through $\tr_k$ cannot
separate two categories agreeing to level $k$, nor a converse realizing every genuine level-$k$
distinction. Note Theorem~\ref{thm:coskel}: for \emph{isomorphism}-invariants the truncation grade is
genuine only up to level two, so any such program must grade \emph{conditions} (or assembly
obstructions), not fixed-category truncations, above that line.

\paragraph{The open question, honestly.}
Does cohomological \emph{degree} stratify \emph{detection} power above the floor? \textbf{No, as
naively stated}, and this is the heart of the correction. Theorem~\ref{thm:degreeone} exhibits a
degree-one invariant (over $\F_2$) that detects composition, so ``nothing below degree two detects
composition'' is false; and Theorem~\ref{thm:sep} shows the degree-one and degree-two anchors are
independent, hence not ordered. What \emph{is} true is the condition-level strictness of
Theorem~\ref{thm:assembly} (composition exists at $\Nerve_2$, is associative at $\Nerve_3$) and the
coefficient-test of Remark~\ref{rem:coefftest} (nerve level is coefficient-free; degree is not). The
grade that stratifies is nerve level; whether degree does anything beyond shadowing it is open.

\paragraph{The sharpened applications story.}
As a grant-narrative frame the ladder is immediately useful, and it sharpens rather than replaces the
earlier picture. Compositional failure is not one phenomenon but several, cleanly separated by the
nerve level of the conserved quantity:
\begin{center}
\begin{tabular}{@{}llll@{}}
\toprule
Datum lives\ldots & Nerve level & Failure class & Economic instance \\
\midrule
on the \textbf{graph} (out-degrees) & $1$ & invisible to $\partial$ & tangent/shape geometry \\
on the \textbf{nodes} (transported) & $2$ & $[\theta_R]\in H^1$ --- route discrepancy & supply-chain inventory \\
in the \textbf{assembly} (weld) & $3$ & $[\omega]\in H^2$ --- handoff order & blockchain re-entrancy \\
\bottomrule
\end{tabular}
\end{center}
In our stylized transport model --- stock carried by translations, not the flow-conservation and
capacity constraints of operations-research inventory accounting --- supply chains give an instance of
``failure $=$ nonzero $H^1$''. The mechanism is classical: a route discrepancy is the obstruction to a
node-potential, exactly Kirchhoff's voltage law, and the same coboundary condition governs
arbitrage-freeness and the ranking inconsistencies measured by HodgeRank \cite{HodgeRank} and by
sheaf-theoretic data fusion \cite{Ghrist}. We therefore do \emph{not} claim supply chains as the
``first'' such instance --- that would need a literature search we have not run --- only as a clean one
in this programme, standing beside re-entrancy's $H^2$ at a different nerve level.

\paragraph{Acknowledgement and honesty.}
Our port of Lanfranchi's theorem \cite{Lanfranchi2026} and our supply-chain conjecture tested in
\cite{MacBethSupplyChain} are both sharpened here, not debunked: the first because its home is a
Cartesian tangent category and the container derivative's is not, the second because inventory is a
node-level degree-one datum and not the weld. The parallelizability definition the port leans on is due
to Cruttwell, Ikonicoff, Lemay and Van Der Linden, as Lanfranchi records; the port itself is ours. The
nerve-level program is a heuristic organizing proved facts; its general form is open, and the slogan it
replaces --- ``composition is a degree-two datum'' --- is retired. Every level assignment in this note
tracks a datum we have computed, and the corrected slogan, ``composition is a nerve-level-two datum,
first seen at cohomological degree one,'' is underwritten by Lemma~\ref{lem:regrade} and
Theorems~\ref{thm:lvl2}, \ref{thm:degreeone} and \ref{thm:thin}.

\begin{thebibliography}{99}

\bibitem{AAGM2003}
M.~Abbott, T.~Altenkirch, N.~Ghani, C.~McBride,
\emph{Derivatives of containers},
Typed Lambda Calculi and Applications (TLCA 2003), LNCS 2701, Springer, 2003, 16--30.

\bibitem{AhmanUustalu2016}
D.~Ahman, T.~Uustalu,
\emph{Directed containers as categories},
Proc.\ MSFP 2016, EPTCS 207, 2016, 89--98; arXiv:1604.01187.

\bibitem{GambinoKock2013}
N.~Gambino, J.~Kock,
\emph{Polynomial functors and polynomial monads},
Math.\ Proc.\ Cambridge Philos.\ Soc.\ \textbf{154} (2013), no.~1, 153--192; arXiv:0906.4931.

\bibitem{Lanfranchi2026}
M.~Lanfranchi,
\emph{The Lie group--Lie algebra correspondence in tangent categories},
preprint, arXiv:2609.03449v1 (2026).
(Parallelizable objects, Def.~4.18, attributed there to
Cruttwell, Ikonicoff, Lemay and Van Der Linden.)

\bibitem{HodgeRank}
X.~Jiang, L.-H.~Lim, Y.~Yao, Y.~Ye,
\emph{Statistical ranking and combinatorial Hodge theory},
Math.\ Program.\ \textbf{127} (2011), no.~1, 203--244.

\bibitem{Ghrist}
R.~Ghrist,
\emph{Elementary Applied Topology},
Createspace, 2014.

\bibitem{MacBethNerve}
MacBeth,
\emph{The nerve-level tower: composition is a nerve-level-two datum (cohomological degree one)}
(Kodamai internal proof, \texttt{2026-10-08-nerve-level-tower}; registry node
\texttt{composition-nerve-level}, proved).

\bibitem{MacBethParallel}
MacBeth,
\emph{Parallelizability does not port to directed containers: a tangent dividing line, and the
out-degree invariant blind to composition}
(Kodamai internal note, 2026; content commit \texttt{e154282}).

\bibitem{MacBethSupplyChain}
MacBeth,
\emph{Supply-chain inventory consistency is a degree-one class, and is independent of the
Zappa--Sz\'ep weld obstruction}
(Kodamai internal note, 2026; content commit \texttt{6f1fbd2}).

\bibitem{MacBethDelta}
MacBeth,
\emph{The $\delta$-left-translation trivialization of the container tangent object sees groupoids}
(Kodamai internal proof, \texttt{2026-10-06-delta-left-translation-trivialization}).

\bibitem{MacBethThin}
MacBeth,
\emph{The constant-presheaf refinement: $\delta$-trivialization with trivial monodromy sees thin
groupoids}
(Kodamai internal proof, \texttt{2026-10-07-constant-presheaf-thin-groupoid}).

\bibitem{MacBethZS}
MacBeth,
\emph{The Zappa--Sz\'ep closure obstruction is a Baues--Wirsching $H^2$ class}
(Kodamai internal note, 2026).

\bibitem{MacBethReentrancy}
MacBeth,
\emph{A machine-checked re-entrancy obstruction: $[\omega]=\varepsilon\in H^2\cong\Z/2$}
(Kodamai internal note, Lean~4 certified, 2026).

\bibitem{MacBethCartan}
MacBeth,
\emph{The container derivative has no fiberwise negation: a boundary for the Cartan calculus of
polynomial functors}
(Kodamai internal note, 2026).

\end{thebibliography}

\end{document}
